%% Comment 3: 3.	Contribution
%Once the ingredients of the policy are clarified, the paper’s contribution should be expanded, because at present it does not fully engage with the most recent literature. Beyond the large body of work on police size and crime, a few recent papers have examined how the allocation of police resources across space and time affects police effectiveness (among others, Adda et al. (2014), Blanes i Vidal and Mastrobuoni (2018), Mastrobuoni (2019, 2020)). Furthermore, it should also connect more explicitly with the growing body of work on the impact of policing on community-related outcomes such as trust, perceptions, and willingness to cooperate (e.g. Blanco et al. (2013), Ang et al. (2025), Facchetti (2025), Gonçalves et al. (2025)).

%I appreciate the effort to relate to work that inform how to organize state presence in space. However, to make this a convincing contribution, the authors would need to discuss explicitly whether centralizing policing over a larger jurisdiction improves the allocation of resources compared to spreading them, and I am not sure if they can currently do it. A cost-benefit analysis, or even a marginal value of public funds (MVPF) estimate, could help.


Thank you for these suggestions. We have read all the cited studies carefully and agree that they are highly relevant and should have been included in our discussion. They are now integrated into the motivation of the paper.

\citet{Adda_etal_2014}, \citet{Blanes_Mastrobuoni_2018}, \citet{MASTROBUONI2019}, \citet{facchetti2025}, and \citet{Mastrobuoni_2020} were particularly helpful in sharpening our understanding of how the allocation of police resources and efforts across space, time and crimes affects policing outcomes and public safety, the mechanisms driving these effects, and of the role that technology can play in improving police effectiveness. We now engage more explicitly with this literature when discussing the contributions of our results.

\citet{Blanco_2025} and \citet{ang_etal_2025} helped us better motivate the importance of community trust in the police, its drivers, and how it shapes policing outcomes, and \citet{Goncalves_etal_2025} documented police misconduct. We now connect our findings more explicitly to this growing literature, highlighting how the intervention we examine can increase trust in the police and the beneficial secondary effects of it. 


Below are the paragraphs in the updated introduction where we have integrated the relevant literature and highlighted the contributions of our results: 

\begin{adjustwidth}{2em}{0em}

Our findings contribute to two related strands of the literature on police organization and public safety. First, they provide some of the first causal evidence on the effects of beat policing on crime and security perceptions. Beat policing has become one of the most widely adopted policing models globally \citep{NASEM2018,PCSP2014}, yet robust causal evidence on its effectiveness remains scarce. We show that this model can generate substantial reductions in crime victimization, reported crime, and insecurity perceptions.

These results also speak to recent work on other policing models. Hot-spot policing reallocates police resources toward high-crime locations and has been shown to reduce crime in targeted areas, though often with displacement to nearby places \citep{DiTella, Blattman_Tobon_JEEA, Draca_2011, Weisburd_restat}. Beat policing differs fundamentally. Rather than concentrating police presence in selected locations---where the same officers need not return repeatedly---it provides continuous, territory-wide coverage through permanent officer-to-area assignment. This distinction matters in our setting, where many quadrants were not crime hot spots and the reform generated municipality-wide declines in crime and insecurity perceptions, with no evidence of spillovers to neighboring municipalities. Community policing, another model that has received growing attention in the economics literature, emphasizes citizen participation in defining policing priorities and, in some cases, in implementation itself. While it has been shown to increase trust in police in developed countries, evidence from low- and middle-income settings points to more limited effects on crime and trust \citep{blair2021, Tobon}. Beat policing also differs from that approach. By permanently assigning officers to stable territorial units, it strengthens local knowledge and brings police closer to communities without requiring direct citizen participation in strategy or operations. Our findings show that this model can reduce crime and insecurity perceptions while increasing trust in police institutions in a middle-income country.

A second contribution is to the broader literature showing that specific features of police organization matter for police effectiveness and public security outcomes. Research has shown that faster police response times improve crime clearance \citep{Blanes_Kirchmaier}, that the effectiveness of police presence depends on how this is organized and deployed rather than simply on their intensity \citep{facchetti2025, MASTROBUONI2019, Blanes_Mastrobuoni_2018}, that shifts in police prioritization across crime types shape criminal behavior \citep{Adda_etal_2014}, and that better information systems and technology can improve police performance \citep{Mastrobuoni_2020}. Earlier work also found that areas covered by smaller, more localized police departments outperform, in terms of security, those covered by larger centralized ones, hypothesizing that this might be related, among other factors, to closer citizen-police engagement \citep{ostrom1973, ostrom1978}. We add to this literature by studying a policing model that combines several of these margins within a single organizational reform. By assigning officers permanently to smaller territorial units rather than rotating them across larger ones, beat policing is designed to increase police visibility, strengthen knowledge of local conditions, and improve the capacity of officers to adapt their practices to local crime patterns. In line with these mechanisms, we show that the adoption of \emph{Plan Cuadrante} increased perceived police presence in neighborhoods and raised police effectiveness, as arrests increased even while victimization and reported crime declined. These organizational improvements translated into meaningful gains in public safety and security perceptions.

A third contribution concerns the literature on police resources, which documents elasticities of crime with respect to officers and expenditure ranging from $-0.1$ to $-2$, with larger effects in developed settings and for violent crimes \citep{Chalfin_JEL, Levitt_1997, DiTella}. We complement this body of work by showing that the quantity of resources is not the only margin that matters. How those resources are deployed matters as well. \emph{Plan Cuadrante} improved public safety outcomes even in municipalities that experienced no increase in police resources relative to controls, suggesting that reorganizing deployment can raise public safety and police effectiveness without expanding the resource base.

We also contribute to the literature on police practices, institutional trust, and citizen cooperation. A persistent challenge is that security perceptions often diverge from actual crime trends \citep{ajzenman2022}, limiting the scope for improving both simultaneously. Rising crime erodes institutional trust \citep{Blanco_2025}, while police misconduct undermines cooperation \citep{ang_etal_2025}.\footnote{Racial disparities in police behavior are also well documented; see, e.g., \citet{Goncalves_etal_2025}, who develop a method to correct for selection bias in estimates of racial disparities in the use of force.} Our results show that beat policing can improve both dimensions---security and security perceptions---by bringing officers into sustained engagement with local communities through permanent territorial assignment. We document not only increases in police trust and evaluations of police performance, but also substantial reductions in private security investment, indicating that households revised their threat assessments in response to the program. 

Finally, our findings speak to a broader question about how organizational design shapes public-service effectiveness. A growing literature shows that the spatial organization of service provision matters. In policing, we show that a model organizing officers around fixed territorial units---allowing them to build local knowledge, coordinate with local actors, and adapt practices accordingly---can substantially improve both crime control and institutional legitimacy. This aligns with evidence from health and education showing that bringing providers closer to the populations they serve can improve outcomes \citep{besley2003centralized, Bardhan, Bloom_2015, Basurto}. In a rigorously identified setting, we show that this type of organizational proximity generates meaningful gains in police effectiveness and citizen security.

\end{adjustwidth}


Regarding the connection of our results with the literature on state service provision organization, we agree that addressing whether centralizing policing over a larger jurisdiction—or, alternatively, spreading officers across smaller, fixed territorial units—improves the allocation of police resources is central to the contribution of the paper. Our findings speak directly to this question. The organizational logic of \emph{Plan Cuadrante} is precisely one of territorial or local specialization: rather than maintaining a centralized dispatch model in which officers patrol in rotation large, undifferentiated areas, the reform assigns officers permanently to small, fixed quadrants, enabling them to build local knowledge and adapt their practices to local crime patterns. Municipalities receiving large increases in police officers and patrol vehicles experienced improvements in victimization similar to those in municipalities receiving resource increases similar to control, indicating that the scale of resource infusion alone does not drive effectiveness. This result supports the view that the organizational model---with permanent assignment to small territorial  areas--- is a key driver of effectiveness. While our design does not allow a direct budget-constrained comparison between a centralized rotating model and a beat policing model, these findings suggest that spreading officers across smaller fixed units, as \emph{Plan Cuadrante} does, generates efficiency gains in resource allocation.

Regarding the request for a cost-benefit analysis or an MVPF estimate, we agree that this would be valuable in principle, but we do not think we can credibly implement it with the data currently available. On the cost side, the budget documents we were able to access (including Budget Office reports) do not cleanly separate baseline police expenditure that would have occurred in the absence of \emph{Plan Cuadrante} from the truly incremental expenditure required to implement and operate the reform. On the benefit side, we also lack a credible context-specific willingness-to-pay measure that would allow us to monetize the welfare gains from the observed reductions in victimization and insecurity. These limitations rule out a credible cost-benefit or MVPF calculation.

That said, an important feature of our setting is that \emph{Plan Cuadrante} is fundamentally an organizational reform of police deployment. Many municipalities that adopted the program experienced resource increases similar to those observed in control municipalities, suggesting that a substantial part of those resources would likely have been deployed even without the reform. In this sense, the gains associated with the beat-policing component are largely linked to how existing police capacity is organized and deployed across territory, rather than to large additional spending above the cost of maintaining a police force. We therefore interpret our findings as evidence that organizational redesign can generate sizeable improvements in public safety even when incremental budgetary costs are limited.


